\clearpage
\phantomsection
\addcontentsline{toc}{section}{Заключение}

\section*{Заключение}

\hspace*{2.5em}В настоящия дисертационен труд бе постигната поставената основна цел – разработване на нов симулатор за ОМ, функциониращ в онлайн среда, който съществено улеснява и ускорява работата с ОМ. Реализираното решение подпомага ключови дейности като създаване на модели, последващо редактиране (включително промени по функции и визуализация) и ефективно споделяне на мрежи. При разработката са използвани съвременни и широко утвърдени технологии, съобразени с изискванията за надеждност, разширяемост и удобство при работа през уеб среда.

Наред със софтуерния резултат, в рамките на дисертацията са получени и нови приноси, свързани с теорията и приложението на ОМ. Разработени бяха нови алгоритми, насочени към подпомагане на моделирането и представянето на ОМ. Първо, предложен е алгоритъм за автоматично изчертаване на ОМ по нейното описание, който съкращава времето за получаване на коректна и прегледна визуализация и ограничава необходимостта от ръчни корекции. Второ, разработен е алгоритъм за експортиране на мрежата във формат \textit{TeX}, така че след извършени редакции и визуални настройки в симулатора, моделът да може директно да бъде включван в научни публикации и други академични материали.

Допълнително бяха създадени нови модели на ОМ в различни приложни области, които демонстрират изразителната мощ на апарата и валидират възможностите на разработения симулатор. Сред тях са модели, свързани с представяне на крайни автомати, модел на математическа игра, както и модели за индустриални процеси, включително сценарии, свързани с производството на газови и нефтени продукти. Всички предложени модели са реализирани и изследвани чрез симулации с разработения в дисертацията симулатор, което осигурява практическа проверка на функционалността и приложимостта на системата.

\paragraph{Насоки за бъдещо развитие.}
\mbox{}\\
Разработеният симулатор дава стабилна основа за надграждане – както откъм нови примерни приложения, така и откъм разширяване на самия ОМ апарат, който системата поддържа. Наред с това, съществуват и редица перспективи за по-нататъшно усъвършенстване на софтуера като програмна система – по отношение на функционалност, производителност, надеждност и сигурност. В този смисъл, като най-естествени следващи стъпки се открояват:

\begin{itemize}
    \item \textbf{Нови ОМ модели и приложения, които досега са разглеждани основно теоретично.}\\
    Практически полезна посока е постепенното разширяване на библиотеката с модели, които са подробно описани в литературата, но рядко са налични в една обща софтуерна среда за създаване, редакция и симулация. Тук попадат модели за представяне и симулация на невронни мрежи чрез ОМ (вкл. с приложения в диагностиката на електронни схеми), моделиране на процеса по извличане на знания чрез ОМ (вкл. \textit{Data Mining}), както и модели за железопътен транспорт, които позволяват анализ и сравнение на различни сценарии и политики. \cite{Sotirov06,Bureva14,Popov20,Boyukov23}
	\newpage
    \item \textbf{Разширение на приложението за поддръжка на други модификации на ОМ.}\\
    Друга перспективна линия е системата да бъде разширена така, че да работи не само със стандартния ОМ формализъм, а и с негови модификации – например \emph{ОМ с характеристики на позициите} и \emph{интуиционистки размити ОМ}. На практика това би изисквало допълнение на средствата за описание на модела (структура, параметри и специфични атрибути), както и адаптиране на симулационния механизъм и визуализацията, така че потребителят да може да моделира и изследва тези разширения по същия удобен начин, както при стандартните ОМ. \cite{Andonov15,Atanassova13}

    \item \textbf{Разширяване на функционалностите на софтуера.}\\
    Като естествено продължение на настоящата разработка може да се предвиди добавяне на нови функционалности, насочени към по-детайлно изследване и по-удобна работа с моделите. Такива биха могли да бъдат допълнителни средства за анализ на симулационните резултати, по-богати възможности за експортиране и документиране на модели, автоматизирани проверки за коректност на входните описания, както и по-гъвкави механизми за конфигуриране и проследяване на изпълнението. Подобни разширения биха повишили както практическата приложимост на системата, така и нейната полезност за изследователски и обучителни цели.

    \item \textbf{Стрес тестване и повишаване на производителността.}\\
    С оглед на бъдещото използване на симулатора с по-големи и по-сложни модели, важна насока за развитие е провеждането на по-задълбочено стрес тестване при различни натоварвания. Това включва изследване на поведението на системата при голям брой преходи, позиции, токени и последователни симулационни стъпки, както и анализ на използваните изчислителни ресурси. На тази основа могат да бъдат предприети оптимизации в алгоритмите, обработката на данните и визуализацията, така че да се подобрят бързодействието, мащабируемостта и общата ефективност на приложението.

    \item \textbf{Повишаване на надеждността и защитата на системата.}\\
    Доколкото \textit{OnlineGN} е уеб-базирано приложение, съществена бъдеща задача е по-нататъшното укрепване на неговата надеждност и сигурност. Това включва както по-строги проверки и валидация на входните данни, така и защита срещу типични атаки към уеб приложения, свързани например с подаване на некоректни или злонамерени входове, претоварване на услугата или опити за неоторизиран достъп. Подобрения в тази насока биха допринесли за по-стабилна експлоатация на системата и биха създали по-добри предпоставки за нейното използване в по-широка изследователска и приложна среда.
\end{itemize}